\documentclass[11pt]{letter}
\usepackage[margin=1in]{geometry}
\usepackage[T1]{fontenc}
\usepackage{lmodern}
\usepackage{hyperref}

\signature{Meng Li, on behalf of Jie Liu, Shiyu Yan, and Xiaohua Yang\\
School of Computing, University of South China\\
Hengyang 421001, China\\
\href{mailto:mlemon@usc.edu.cn}{mlemon@usc.edu.cn}}

\begin{document}

\begin{letter}{Professor Mauro Pezze\\
Editor-in-Chief\\
ACM Transactions on Software Engineering and Methodology}

\opening{Dear Professor Pezze,}

We are pleased to submit our manuscript ``NOETHER: Constructive Metamorphic Pattern Identification from Operator Algebras and a Falsifiable Invariance-Blindness Theorem'' for consideration in ACM TOSEM under the \textbf{Fast Impact} track.

The manuscript has been shortened so that the main paper body fits the TOSEM Fast Impact 45-page limit in the official \texttt{acmsmall,screen,anonymous,review} format. References and online-only appendix/supplementary material are kept outside the main-body count. Self-contained implementation details, detailed derivation traces, and secondary tables have been migrated to appendices and supplementary material (S1--S12).

Metamorphic testing depends on metamorphic relations, yet MR identification remains experience- or search-driven, with no structural account of where MR patterns come from, when a pattern set is closed, or how it transfers across program families. This paper addresses that origin--closure--transferability gap through five contributions:

\begin{enumerate}
\item A layered, constructive framework (NOETHER) that derives MetaPatterns from the operator-algebraic structure of a program family by the mechanical CONSTRUCT-MP algorithm.
\item A no-drop closure invariant and polynomial-time constructibility result over the algebra-induced MR space.
\item A falsifiable Invariance-Blindness Theorem characterising exactly which faults symmetry and self-adjoint MetaPattern MRs can and cannot detect within the stated fault class.
\item A negative PWR reactor-physics instantiation showing that absolute completeness is false and identifying five independent obstructions.
\item A structural-transferability and evidence protocol instantiated on Boltzmann reactor physics, equivariant machine learning, and relational query optimisers.
\end{enumerate}

A preprint of this work is available on arXiv (arXiv:2605.17390, \url{https://arxiv.org/abs/2605.17390}). The submitted manuscript is the double-blind version; the arXiv identifier and the preprint label have been removed from the anonymized PDF.

A separate line of work by a subset of the authors addresses the orthogonal problem of selecting a minimum complete subset from a given MR pool. That work takes an MR set as input and minimises its cardinality under a fixed fault model. The present paper concerns the upstream question of where MRs and MetaPatterns come from: their constructive derivation from a program-induced operator algebra and algebraic closure under the \texttt{Translate} operator. The two share no theorem or empirical claim.

In line with ACM's Policy on Authorship, large language models appear in this work in two roles. As instruments of the study, they are part of the reported LLM baseline, LRCA second-rater protocol, and mutant-equivalence adjudication; the prompts and raw outputs are released with the replication package. As authoring assistance, the authors used a large-language-model assistant for code scaffolding, reference formatting, and language editing. All research design, theorems, proofs, experimental procedures, numerical results, and claims were produced and verified by the authors, who take full responsibility for the manuscript.

The authors declare no competing interests.

\closing{Sincerely,}

\end{letter}

\end{document}
